% variants: plain tag
% Tests: float.sty [H] floats, nested [H] float inside a minipage, \restylefloat, \newfloat.
\documentclass{article}
\usepackage{float}
\usepackage{caption}
\floatstyle{boxed}
\restylefloat{table}
\floatstyle{ruled}
\newfloat{program}{tbp}{lop}
\floatname{program}{Program}
\captionsetup[figure]{labelfont=bf}
\begin{document}
\listoffigures
\listoftables
\listof{program}{List of Programs}
\section{Here floats}
\begin{figure}[H]
\centering\rule{3cm}{1cm}
\caption{An H figure}\label{fig:h}
\end{figure}
\begin{table}[H]
\caption{A boxed H table}\label{tab:h}
\centering\rule{3cm}{1cm}
\end{table}
\begin{program}[H]
\caption{A ruled program}\label{prog:h}
\texttt{x := 1}
\end{program}
\noindent\begin{minipage}{.45\linewidth}
\begin{figure}[H]
\centering\rule{2cm}{1cm}
\caption{Nested H figure in a minipage}\label{fig:nested}
\end{figure}
\end{minipage}\hfill
\begin{minipage}{.45\linewidth}
\begin{table}[H]
\caption{Nested H table in a minipage}\label{tab:nested}
\centering\rule{2cm}{1cm}
\end{table}
\end{minipage}
See \ref{fig:h}, \ref{tab:h}, \ref{prog:h}, \ref{fig:nested}, \ref{tab:nested}.
\end{document}
